\section{Exterior tensors at a prescribed placement}\label{sec:exterior}

The star centered at the cone tip supplies a tensor by sending its
indices separately to their sites. To detect the saving obtained by
using a separate supplier on each side of a balanced cut, we need a
tensor whose price at the original sites is already close to the cost
of the star. We first construct such a tensor for a prescribed
placement. In one copy, its nonzero entries are maximal minors of a
matrix selected for those sites and their spectral price.
Section~\ref{sec:packing} will group copies by site and retain certain
row letters to obtain a weighted tensor by restriction from one
explicit binary tensor.

Let $Y$ be a complete Hadamard cone with tip $0$ and homogeneous spectral
prices over an algebraically closed field $K$. We use natural logarithms.
For a $b\times N$ matrix $C$, consider $N$ inputs with two states,
empty and occupied. The coefficient on the set of occupied inputs $I$
is the ordered minor $\det C_I$ when $|I|=b$, and zero otherwise.
For example, when the columns are $(1,t_i)^{\mathsf T}$ and $b=2$,
the coefficient on $\{i,j\}$, $i<j$, is $t_j-t_i$. The exterior
product of the occupied columns produces these coefficients and lets
us keep track of the number of times each input is used in many copies.

More generally, let $A_i\subseteq E$ be subspaces spanning a
finite-dimensional space $E$. For $\ell\ge1$, wedge the inputs from
$\bigwedge^*(A_i\otimes K^\ell)$ in a fixed order, take the component
of top degree, and apply a volume form on $E\otimes K^\ell$. Write
$\Omega_\ell((A_i))$ for this multilinear form. Its \emph{sector}
$\Omega_{\ell,p}$ fixes the exterior degree $p_i$ at each input, with
$\sum_i p_i=\ell\dim E$. A \emph{supplier} $(A_i,y_i,\theta_i)$ puts
this input at the site $y_i$ and credits it $\theta_i$ units per degree.
The value of a list of suppliers is the largest biased sector price per
copy, optimized over the number of copies,
\begin{equation}\label{ext:value}
 \lambda((A_i,y_i,\theta_i)_i)
 =\sup_{\ell\ge1}\frac1\ell
   \max_p\left\{H(\Omega_{\ell,p};y)+\sum_i\theta_i p_i\right\}.
\end{equation}
Linear maps into $E$ may replace the inclusions $A_i\subseteq E$.
Factoring a map through its image and choosing a linear section gives
restrictions in both directions, so the sector prices agree. Taking
one map $K\to E$ per column of $C$ gives the form $\Omega_\ell(C)$
and the value $\lambda(C;y,\theta)$.

A column of degree $p$ in $\ell$ copies has local dimension
$\binom{\ell}{p}$. Sending this input from the tip to a site at
distance $r$ costs at most $r\log\binom{\ell}{p}$. If $h$ denotes
binary entropy in nats, its biased price per copy is bounded by
$rh(p/\ell)+\theta p/\ell$. Maximizing the latter over $p/\ell\in[0,1]$
gives the budget for a single column,
\[
 A(r,t)=r\log(1+e^{t/r})\quad(r>0),\qquad
 A(0,t)=\max(0,t),
\]
where $t$ is its bias. Summing these budgets gives the star upper bound
$\lambda(C;y,\theta)\le\sum_i A(\|y_i\|,\theta_i)$.
We use Vandermonde matrices with columns
$(1,t_i,\ldots,t_i^{b-1})^{\mathsf T}$ for distinct $t_i\in K$.
Any $b$ of these columns form a basis of $K^b$.

\begin{theorem}[Exterior near-star estimate]\leavevmode\phantomsection\label{ext:star}
Let $x_1,\ldots,x_d$ be unit sites with uniform Cartan barycenter $0$.
Give each site $q$ column suppliers, with identical lists of biases
\[
 \theta_j=-M+(2j-1)M/q,\qquad 1\le j\le q.
\]
Suppose $N=dq$ is even and $M\ge100\log N$. Write $y,\theta$ for
the resulting lists of $N$ sites and biases. There exists an
$(N/2)\times N$ Vandermonde matrix $C$ satisfying
\[
 0\le \sum_i A(1,\theta_i)-\lambda(C;y,\theta)
       \le C_0 M(1+N/q^2)\log N
\]
with an absolute constant $C_0$.
\end{theorem}

To see what the lower bound asks of the columns, give the tip an
additional supplier of the whole space $E$ and vary its bias $s$.
As $s$ increases, it becomes advantageous to take more degrees at
the tip. The marginal value of $s$ measures those degrees per copy.
Subtracting it from $b=\dim E$ gives the \emph{threshold rank} $r_s$
still supplied by the original columns. Equation~\eqref{ext:recover}
will recover the original price by integrating $r_s$.

A single column at a unit site has threshold rank
$p_i(s)=(1+e^{s-\theta_i})^{-1}$, and
$\int_0^\infty p_i(s)\,ds=A(1,\theta_i)$. To recover the star budget
we would like $r_s$ to be close to $\sum_i p_i(s)$ above zero, where
the columns must contribute independently, and close to $b$ below
zero, where they must fill $E$. For an ordinary subset of Vandermonde
columns these two requirements meet in the identity
$\rank(C_I)=\min(b,|I|)$.

For priced columns we can construct an event of dimension $p_i(s)$ for
each column. Intersection dimensions give a probability law for each
\emph{pair} of events. These laws need not extend to a joint law for all
the events, so in general we cannot sample a subset $I$ and count its
members. Section~\ref{sec:flags} inserts event indicators one at a
time and shows how the resulting count bounds the Vandermonde rank from
below.
Section~\ref{sec:achieve} uses the centered sites and the evenly spaced
biases to control the second moment of this count over insertion orders.
This gives one choice of columns with small integrated rank loss.

\subsection{Exterior price rules}

We record the properties of the exterior value that underlie the
threshold rank. First, the supremum in~\eqref{ext:value} is a limit.
Split the copy indices into blocks of lengths $\ell_1,\ell_2$ and fix
a maximizing sector in each block. Wedge the two inputs locally at
every supplier. The sector of the combined system then restricts to
the product of the two sectors. The sign for moving the blocks past
one another is constant on the fixed degrees. Hence the unnormalized
maxima are superadditive. They are bounded above by $O(\ell)$ by
transport, and spanning gives a nonzero sector. Fekete's lemma now
yields the limit. Reordering suppliers also changes each sector only
by a scalar sign. For the full exterior form this sign can depend on
the degrees of the inputs.

Enlarging or adding suppliers gives restrictions onto the original
sectors. Increasing a bias can only increase the value, and shifting all biases
by $a$ adds $a\dim E$ because the total degree is fixed. Joint convexity
and positive homogeneity in the labels follow from the corresponding
properties of price. Continuity follows from transport and the bound
$\ell\dim A_i\log 2$ on the logarithm of each local dimension. A zero
subspace contributes nothing.

Two suppliers at the same site and bias may be merged by replacing their
subspaces with their sum. Precomposing the merged form with local
wedge multiplication gives the form on each split sector. For the reverse inequality, fix a
merged degree and take the direct sum of the wedge maps over all degree
splits. This map is onto because the two subspaces span their sum. A
linear section expresses the merged tensor as an ordinary sum of
restrictions, one for each of at most $\ell\dim E+1$ splits. Their
bias credits agree. Lemma~\ref{price:rules} bounds the extra price by
$O(\log(\ell+1))$, which vanishes per copy. Direct sums of systems in
complementary summands of $E$ add values, since their nonzero sectors
factor.

Finally, consider lists $L_1,L_2$ with a spanning common sublist.
We prove the diminishing-increments inequality
\[
 \lambda(L_1\cup L_2)+\lambda(L_1\cap L_2)
 \le \lambda(L_1)+\lambda(L_2).
\]
Place the lists in separate summands of $E\oplus E$ and merge common
suppliers. For $t\ne0$ apply the change of coordinates
$(u,v)\mapsto(u+v,tv)$. Each common supplier still supplies
$A_i\oplus A_i$ and can keep its original basis. Use the transformed
bases, polynomial in $t$, for all exclusive suppliers. At $t=0$ they
supply subspaces of the first summand. Each limiting sector is thus
a specialization of isomorphic sectors, so its price cannot increase
by Lemma~\ref{price:rules}. Splitting the common suppliers again gives
the union in the first summand and the intersection in the second.
Additivity now gives the displayed inequality. For lists with repeated
members, use named occurrences to form the union and intersection.
